% !TITLE 燕歌行其一
% !AUTHOR 曹丕
% !DYNASTY 魏晋
% !GENRE 七言古诗
% !ORDER 5

\begin{poem}
秋风萧瑟天气凉，草木摇落\textsuperscript{1}露为霜。
群燕辞归鹄南翔，念君客游思断肠。
慊慊\textsuperscript{2}思归恋故乡，君何淹留\textsuperscript{3}寄他方？
贱妾茕茕\textsuperscript{4}守空房，忧来思君不敢忘，不觉泪下沾衣裳。
援琴鸣弦发清商\textsuperscript{5}，短歌微吟不能长。
明月皎皎照我床，星汉西流夜未央\textsuperscript{6}。
牵牛织女遥相望，尔独何辜限河梁\textsuperscript{7}。
\end{poem}

\begin{annotations}
\notetext{1}{摇落：凋残零落。常用来形容秋天草木枯萎的景象。}
\notetext{2}{慊慊：慊慊（qiàn qiàn），因思念、不满或失落而感到心里空虚忧伤的样子。}
\notetext{3}{淹留：长期滞留、逗留。}
\notetext{4}{茕茕：茕茕（qióng qióng），孤单无依的样子。}
\notetext{5}{清商：古代的乐曲名。声调凄清哀怨，适合表达悲伤的心情。}
\notetext{6}{未央：没有尽头，指夜还很深。央，尽、结束。}
\notetext{7}{河梁：河上的桥。这里指离别时走过的桥，也比喻阻碍亲人相聚的地方。}
\end{annotations}

\begin{appreciation}
《燕歌行》是中国文学史上现存最早、最完整的一首七言诗之一，作者曹丕——曹操的儿子，魏国的皇帝。一个皇帝写女子秋夜思念丈夫，写得比闺中女儿还细。

秋风萧瑟天气凉，草木摇落露为霜——秋风起，草木凋零，露水结霜，先铺开一个秋天。群燕辞归鹄南翔，念君客游思断肠——燕子大雁都往南飞了，想着你在外乡，想得肝肠寸断。鸟都知道赶在冬天前回去，人却不回来。慊慊思归恋故乡，君何淹留寄他方？——她明白丈夫也想家，可你想家，怎么还留在远方？这一问里有理解，也有委屈。

贱妾茕茕守空房，忧来思君不敢忘，不觉泪下沾衣裳——孤零零守着空房，忧愁起来就想你，泪不知不觉打湿衣裳。"不敢忘"三个字最好：不是不能忘，是不敢忘——好像忘了你是一种罪过。援琴鸣弦发清商，短歌微吟不能长——想弹琴解闷，弹的又是哀怨的调子，越弹越闷。明月皎皎照我床，星汉西流夜未央——月光照床，银河西移，夜还深着，她睡不着，躺着看星星。

最后两句，她望见牵牛织女隔河相望，问：你们有什么罪过，要被河水隔在两岸？由自己的离别想到天上的离别，一个人的苦变成了所有人的苦。

还记得《行行重行行》吗？"衣带日已缓"，那里的女子等游子等到消瘦。从古诗十九首到燕歌行，思妇诗一脉相承：都是秋天，都是夜晚，都在等一个不回来的人。句子从五言变七言，那份"等"没有变。

秋天是容易想家的季节。天凉了，自己记得加衣服。这话不是诗里的，是我要说的。
\end{appreciation}
